长江流域是我国淡水渔业最发达、生物多样性最丰富的区域之一，但长期受过度捕捞、工业废水排放、水域塑料污染以及葛洲坝、三峡大坝等水利工程的综合影响，鱼类洄游通道被阻断，渔业资源持续衰退，典型食物链结构发生退化。自2021年1月1日实施长江十年禁渔政策以来，长江流域生态环境已出现显著变化：局部水域水质改善，水草等水生植被逐步恢复；2025年监测数据显示，青、草、鲢、鳙四大家鱼资源量已恢复至禁渔前的约1.8倍，表明禁渔对基础资源和经济鱼类具有显著的修复作用。
